\begin{resumo}[Abstract]
 \begin{otherlanguage*}{english}
The Integrated Academic Activities Management System (SIGAA) is the single access point for University of Brasília (UnB) students to grades, schedules, transcripts, documents and virtual classrooms, with no institutional alternative. Recurring reports of usability problems, a desktop-oriented interface and the existence of community tools that replace system functions motivate an objective assessment of its quality. This paper presents the specification of that assessment, carried out by the Adele Goldstine team following the SQuaRE evaluation process (ISO/IEC 25040) and the ISO/IEC 25010:2023 product quality model. In Phase 1, the purpose, stakeholders, tailored quality model and scope were defined, and an Impact~$\times$~Risk matrix with a feasibility filter selected four characteristics: Interaction Capability (Usability), Functional Suitability, Performance Efficiency and Flexibility (Portability). In Phase 2, the Goal-Question-Metric (GQM) method refined them into four goals, 14 questions with hypotheses and 20 metrics, each with a collection procedure, rating levels and assessment criteria fixed before data collection. Project management is tracked with PSM/CID indicators.

   \vspace{\onelineskip}

   \noindent
   \textbf{Keywords}: software product quality. SQuaRE. GQM. SIGAA.
 \end{otherlanguage*}
\end{resumo}
